% TikZ 图：尖端放电与避雷针
% Day 15：电荷、电场与电场强度
\begin{tikzpicture}[scale=0.95,>=Latex,line cap=round,line join=round]
  \path[use as bounding box] (-4.4,-1.65) rectangle (4.4,2.4);

  % Charged cloud.
  \draw[thick,gray!70] (-3.9,1.85)
    .. controls (-3.6,2.25) and (-3.0,2.20) .. (-2.75,1.88)
    .. controls (-2.35,2.25) and (-1.72,2.10) .. (-1.58,1.76)
    .. controls (-1.05,1.95) and (-0.62,1.66) .. (-0.78,1.32)
    .. controls (-1.55,1.18) and (-2.75,1.20) .. (-3.9,1.35)
    .. controls (-4.15,1.48) and (-4.16,1.70) .. (-3.9,1.85);
  \foreach \x/\y in {-3.35/1.65,-2.70/1.78,-2.05/1.55,-1.35/1.60} {
    \node[blue!70!black,font=\small] at (\x,\y) {$-$};
  }
  \node[font=\small] at (-2.35,2.28) {带电云层};

  % Ground and building.
  \draw[thick] (-4.2,-0.95) -- (4.2,-0.95);
  \foreach \x in {-4.0,-3.55,...,4.0} {
    \draw[gray!60] (\x,-0.95) -- ++(-0.18,-0.12);
  }
  \draw[fill=gray!12,thick] (1.3,-0.95) rectangle (2.75,0.25);
  \draw[fill=gray!18,thick] (1.18,0.25) -- (2.02,0.82) -- (2.88,0.25) -- cycle;
  \node[font=\small] at (2.03,-0.35) {建筑};

  % Lightning rod.
  \draw[very thick] (2.02,0.82) -- (2.02,1.65);
  \draw[fill=black] (2.02,1.65) -- (1.92,1.42) -- (2.12,1.42) -- cycle;
  \draw[very thick] (2.02,0.82) -- (2.02,-1.20);
  \node[font=\scriptsize,right=2pt] at (2.05,-1.18) {接地};

  % Dense positive charges near sharp tip.
  \foreach \x/\y in {1.86/1.47,1.95/1.56,2.05/1.56,2.17/1.47,2.02/1.33} {
    \node[red!80!black,font=\scriptsize] at (\x,\y) {$+$};
  }

  % Field lines around the sharp point.
  \foreach \a in {72,92,112,132} {
    \draw[red!75!black,->,thick] (2.02,1.55) -- ++(\a:0.90);
  }
  \node[red!80!black,font=\small,align=center] at (3.45,1.88)
    {尖端附近\\电场更强};

  % Smooth conductor comparison.
  \draw[fill=gray!18,thick] (-2.10,-0.95) rectangle (-1.25,0.28);
  \draw[fill=gray!18,thick] (-2.10,0.28) arc[start angle=180,end angle=0,radius=0.425] -- cycle;
  \foreach \x/\y in {-1.98/0.48,-1.72/0.62,-1.45/0.48} {
    \node[red!80!black,font=\scriptsize] at (\x,\y) {$+$};
  }
  \draw[red!75!black,->] (-1.95,0.55) -- (-2.35,0.95);
  \draw[red!75!black,->] (-1.68,0.68) -- (-1.68,1.15);
  \draw[red!75!black,->] (-1.42,0.55) -- (-1.02,0.95);
  \node[font=\small,align=center] at (-1.68,-1.35)
    {圆滑表面\\电荷较分散};

\end{tikzpicture}
